\section{El tema central de la infraestructura en 2025: de la clusterización de la inferencia a la infraestructura de RL / Agent}

\subsection{Por qué la infraestructura ocupó por primera vez el centro del escenario}
La apertura del panel planteó el problema con mucha claridad: en 2025 la infraestructura ya no es solo un soporte de fondo, sino la condición previa para que la capacidad de un modelo pueda materializarse. Varios de los invitados mencionaron una sensación compartida: en el pasado, la discusión de la industria se concentraba más en el modelo mismo, la innovación algorítmica y la escala de los datos, pero hacia 2025 empezó a repetirse la conciencia de que ``el mismo modelo, si la infraestructura no funciona, no puede liberar su capacidad bajo una carga real''.

\begin{importantbox}{Por qué la importancia de la infraestructura aumentó de manera notoria en 2025}
Hay al menos tres razones:
\begin{itemize}
  \item la inferencia pasó de experimentos en una sola máquina a producción a gran escala, por lo que el rendimiento, la latencia y la estratificación de recursos deben diseñarse de forma sistemática;
  \item RL, Agent y multi-modality trajeron consigo rollouts más largos, workloads más heterogéneos y cadenas de retroalimentación más complejas;
  \item la industria ya no solo entrena modelos, sino que empieza a entrenar todo un sistema de servicio capaz de interactuar con sistemas externos.
\end{itemize}
\end{importantbox}

Este cambio significa que la infraestructura ya no consiste solo en ``hacer que los operadores corran más rápido'', sino que empieza a determinar el paradigma mismo de la investigación. Por ejemplo, cuando se discute si Agent RL es viable, en el fondo se discute si el motor de rollout, la construcción del entorno, la combinación de entrenamiento e inferencia y la recuperación ante fallas son suficientes para sostener el ciclo cerrado del experimento.

\subsection{La palabra clave del primer semestre: la clusterización de la inferencia}
Zhang Mingxing (张明星) ofreció una división muy nítida: el tema del primer semestre de 2025 fue la clusterización de la inferencia, mientras que el del segundo semestre se inclinó más hacia el scaling de la infraestructura de RL. Esta división resulta muy reveladora, porque muestra que la industria no resuelve todos los problemas al mismo tiempo, sino que primero convierte la inferencia en línea en una base industrializada, y solo después conecta a ella workloads más complejos como RL y Agent.

\textit{``En el primer semestre, en realidad todo se trataba de la clusterización de la inferencia; el tema del segundo semestre fue, básicamente, cómo hacer scaling up / scaling out de la infraestructura de RL.''}

\begin{knowledgebox}{Por qué la clusterización de la inferencia es condición previa de la infraestructura de RL / Agent}
Porque el cuello de botella de muchos workloads nuevos no está en el backprop, sino en el rollout, el sampling y la interacción con el entorno. Dicho de otro modo, si el motor de inferencia no puede funcionar de manera estable, con alto rendimiento y baja fragmentación, entonces el entrenamiento posterior de RL / Agent simplemente no cuenta con un ciclo de retroalimentación suficientemente barato.
\end{knowledgebox}

\subsection{Resumen de la sección}
La infraestructura fue importante en 2025 no porque la ingeniería de sistemas se haya vuelto de repente un tema de moda, sino porque tanto el entrenamiento de modelos como el despliegue de productos se vieron forzados a una etapa en la que ``deben sistematizarse''. La clusterización de la inferencia se convirtió primero en el eje del primer semestre, y después la infraestructura de RL / Agent tomó el relevo, conformando el eje central de la discusión de todo el año.
